\begin{proof}
各 $B_i$ を
\[
B_i = I_{i,1} \times \cdots \times I_{i,d}
\]
と書く．
各座標方向 $k=1,\dots,d$ について，
$I_{i,k}$ に現れる全ての端点を集めて
\[
t_{k,0} < t_{k,1} < \cdots < t_{k,N_k}
\]
と並べる．
これらでできる格子の小半開区間
\[
Q_\alpha
:=
\prod_{k=1}^d [t_{k,\alpha_k-1},t_{k,\alpha_k})
\qquad
(1 \le \alpha_k \le N_k)
\]
を考える．

まず $B_i$ が半開区間である場合を考える．
各座標方向で
$[t_{k,\alpha_k-1},t_{k,\alpha_k})$ は
$I_{i,k}$ に含まれるか，さもなくば交わらないかのどちらかであるから，
各格子小区間 $Q_\alpha$ についても
\[
Q_\alpha \subset B_i
\quad\text{または}\quad
Q_\alpha \cap B_i = \emptyset
\]
のいずれかである．
したがって
\[
A
=
\bigsqcup_{Q_\alpha \subset A} Q_\alpha
\]
と表せるので，$A$ は基本集合として書ける．

一般の有界区間の場合も，端点の開閉の違いは格子の境界上でしか起こらないから，
同じ端点をもつ半開区間に直して上と同じ議論をすればよい．

次に体積和が表示に依らないことを示す．
$A$ の二つの基本集合表示
\[
A = \bigsqcup_{j=1}^m Q_j = \bigsqcup_{\ell=1}^r Q'_\ell
\]
を取る．
$Q_j$ と $Q'_\ell$ に現れる全ての端点を集めて共通の格子分割を作ると，
その各小区間 $S_\nu$ は各 $Q_j$ に対しても各 $Q'_\ell$ に対しても
「含まれる」か「交わらない」かのどちらかである．
ゆえに
\[
A = \bigsqcup_{S_\nu \subset A} S_\nu
\]
であり，しかも各 $Q_j$ も各 $Q'_\ell$ も
これらの $S_\nu$ の互いに素な合併に分解される．
したがって
\[
\sum_{j=1}^m |Q_j|
=
\sum_{S_\nu \subset A} |S_\nu|
=
\sum_{\ell=1}^r |Q'_\ell|
\]
となる．
よって，書き換え後の体積 $\sum_{j=1}^m |Q_j|$ は分割の方法に依らない．
\end{proof}
